\section{Discussion and Conclusion}
\label{sec:discussion}

\subsection{What the Full Study Supports}

The completed study supports a stronger and narrower claim than the pilot alone. At 0.5B, countermodel supervision mainly acts as a faithfulness mechanism. At 7B, the core result is not universal raw-accuracy dominance, but a robust gain in verifier-backed reasoning quality. On the fixed 4k seeded suite, \model{} more than doubles proof-only joint accuracy in-domain and on depth-OOD. As the training budget grows, that reasoning gain persists while the raw-accuracy gap steadily contracts. The tradeoff is therefore \emph{scale-sensitive}, not intrinsic.

\subsection{Why \model{} Helps}

The most important insight is that the benefit does not come primarily from predicting \texttt{Unknown} more often. At 7B, proof-only and \model{} abstain on nearly the same number of gold-unknown examples. The difference is that proof-only abstentions remain generic and unverifiable, whereas \model{} abstentions become explicit missing-support witnesses. In other words, the crucial ingredient is not abstention as a label, but abstention as a structured evidence object. The ablation table reinforces this interpretation: both failure tags and witness text matter, and the full target works best when the two are combined.

\subsection{Scaling Behavior}

Scaling changes the paper's interpretation in an important way. At low budget, \model{} should be read as a method for buying much better faithfulness at a modest raw-accuracy cost. At larger budgets, that reading becomes incomplete. On the full-train fixed depth-5 subset, \model{} slightly surpasses answer-only in raw accuracy, and on several transfer settings it remains competitive while preserving very large joint-accuracy gains. The practical lesson is that countermodel-aware supervision becomes more attractive, not less, as the base model and training budget improve.

\subsection{Limitations}

The paper is still not a ``solved logical reasoning'' result. First, the study uses one benchmark family and one main 7B backbone, so the strongest claims should be restricted to ProofWriter-style reasoning. Second, answer-only systems receive joint accuracy 0 by construction because they do not emit evidence; this is appropriate for a faithfulness paper, but it means that joint accuracy should be interpreted alongside raw accuracy rather than in isolation. Third, the abstention ablations are single-seed measurements, which are enough to localize the effect but not enough to fully characterize variance.

\subsection{Next Steps}

The natural next step is to stress-test the same supervision idea beyond the current setting. The most useful directions are:
\begin{itemize}
\item \textbf{backbone breadth}: repeat the same supervision comparison on additional 7B--14B models to test whether the scaling pattern is architecture-specific,
\item \textbf{benchmark breadth}: evaluate on other verifier-backed reasoning datasets where abstention and refutation matter,
\item \textbf{harder OOD settings}: target the remaining gap on long refutation chains and more adversarial distribution shifts,
\item \textbf{answer/rationale decoupling}: combine a short answer head with a verifier-constrained rationale channel so that raw classification and evidence generation can specialize.
\end{itemize}

\subsection{Conclusion}

We presented \model, a proof-or-countermodel distillation format for logical reasoning. The main result is not that countermodel supervision dominates every raw metric, but that it reliably improves verifier-backed reasoning and abstention faithfulness, and that this advantage survives seeds, scaling, and transfer. Proof-only supervision can learn to output structured traces, but it still fails to make abstention evidence verifiable. \model{} closes that gap, making countermodels a practical supervision target for faithful logical reasoning while leaving broader benchmark transfer and harder long-depth reasoning as the next frontier.
